\documentclass[12pt,a4paper]{article}
\usepackage[utf8]{inputenc}
\usepackage[french]{babel}
\usepackage[margin=2.5cm]{geometry}
\usepackage{enumitem}
\usepackage{amsmath}

\title{Mécanismes de jeu}
\author{}
\date{}

\begin{document}

\maketitle

\section{Loot}
Tirer 1d6 :
\begin{itemize}[label=--]
    \item \textbf{1--2} : choisir 1 item
    \item \textbf{3--4} : choisir 2 items
    \item \textbf{5--6} : tout loot
\end{itemize}

\section{États et afflictions}

\subsection{État enflammé}
$-1$ au début du tour de jeu de la personne/créature dans cet état.

\subsection{État empoisonné}
$-1$ au début du tour de jeu de la personne/créature dans cet état.

\subsection{État étourdi}
Malus de $-2$ aux dés.

\subsection{État assommé}
Passe 1 tour.

\section{Combat}

\subsection{Attaquer}
Lancer 2d10 puis calculer :
\begin{align*}
\text{Total} &= 2\text{d}10 + \text{FOR (mêlée) ou AGI (distance)} + \text{bonus} - \text{RES cible} - \text{malus}
\end{align*}

La cible subit : $\text{Total} - 10$ points de dégâts.

\textbf{Note :} Si le résultat est négatif, c'est l'attaquant qui subit les dégâts (sa RES s'applique également) car on considère qu'il y a riposte.

\subsection{Coup critique}
Après avoir tout calculé, on double les dégâts qui devraient être subis (donc après décompte de la RES, etc.).

\section{Soins et repos}

\subsection{Kit de soin}
Soigne 1d6 points de vie.

\subsection{Repos nocturne}
Tirer 1d6 :
\begin{itemize}[label=--]
    \item \textbf{1--2} : regain de 1 PV
    \item \textbf{3--4} : regain de 2 PV
    \item \textbf{5--6} : regain de 3 PV
\end{itemize}

\section{Événements}

\subsection{Événement nocturne}
Tirer 1d6 :
\begin{itemize}[label=--]
    \item \textbf{1--2} : mauvaise rencontre avec une créature ou, à défaut, un événement négatif
    \item \textbf{3--4} : rien à signaler
    \item \textbf{5--6} : événement positif qui apporte un loot au choix parmi quelques items du biome (éventuellement ceux manqués sur des créatures passées)
\end{itemize}

\section{Situations spéciales}

\subsection{Duel psychique contre Hollow}
Pour échapper au contrôle mental de Hollow, lancer 2d10 de VOL (volonté) :
\begin{itemize}[label=--]
    \item Si Hollow fait plus : contrôle
    \item Sinon : il passe son tour et retourne dans l'ombre
\end{itemize}

\end{document}
